\section{Catastrophic Cancellation and Stable Reformulations}

\begin{remark}[Catastrophic Cancellation]
When large, nearly equal quantities cancel, an error previously
negligible relative to the intermediate values can dominate the
final result. The preceding example loses the small term $a$ in
the first rounded addition.

The same phenomenon can reveal \emph{pre-existing input errors}.
If $\widehat{x},\widehat{y}\in\mathbb F$ approximate $x,y\in\R$,
then even when their subtraction is exact,
\[
 |(\widehat{x}-\widehat{y})-(x-y)|
 \le |\widehat{x}-x|+|\widehat{y}-y|.
\]
For $x\ne y$, dividing by $|x-y|$ gives the relative error;
dividing by $\ulp(x-y)$ gives the ULP error. Both can be large
when $x\approx y$, even if the subtraction itself is exact.
\end{remark}


\begin{example}[Stable Evaluation of $a^2-b^2$]
For $a, b\in\mathbb F$, $a^2\ne b^2$, compare
\[
 u=(a\otimes a)\ominus(b\otimes b), \qquad
 v=(a\ominus b)\otimes(a\oplus b).
\]
Assume that the operations do not overflow or underflow and that the
usual relative-error model applies to all intermediate results.
Writing $|\delta_i|\le\varepsilon$ and retaining terms through first
order gives
\[
\begin{aligned}
 u&=(a^2(1+\delta_1)-b^2(1+\delta_2))(1+\delta_3), \\
 \frac{|u-(a^2-b^2)|}{|a^2-b^2|}
 &\le\left(1+\frac{a^2+b^2}{|a^2-b^2|}\right)\varepsilon
   +O(\varepsilon^2), \\
 v&=(a-b)(a+b)(1+\delta_1)(1+\delta_2)(1+\delta_3), \\
 \frac{|v-(a^2-b^2)|}{|a^2-b^2|}
 &\le3\varepsilon+O(\varepsilon^2).
\end{aligned}
\]
For a normal exact result, the corresponding ULP-error bounds are
$1+(a^2+b^2)/|a^2-b^2|+O(\varepsilon)$ and $3+O(\varepsilon)$,
respectively. The first can be arbitrarily large when $|a|\approx|b|$;
the factored expression avoids subtracting two rounded squares.
\end{example}

\begin{remark}
For $r, r'\in\R$ with
$\operatorname{err}_{\mathrm{ulp}}(r, r')\le1$ and finite $\rnd(r')$,
\[
 \operatorname{err}_{\mathrm{ulp}}(r, \rnd(r'))
 \le\operatorname{err}_{\mathrm{ulp}}(r, r')+\tfrac12.
\]
Indeed, if $\ulp(r')\le\ulp(r)$, the triangle inequality and nearest
rounding suffice. Otherwise $r'$ lies across a power-of-two boundary
from $r$. The one-ULP hypothesis forces $r'$ to round back to that
boundary, closer to $r$ than $r'$ was. The estimate thus holds in both
cases.
\end{remark}

\begin{example}[Stable Quadratic Roots]
For $a\ne0$ and $\Delta=b^2-4ac\ge0$, the usual quadratic formula is
\[
 x_{\pm}=\frac{-b\pm\sqrt\Delta}{2a}.
\]
If $b<0$ and $b^2\gg|4ac|$, evaluating
$(-b-\sqrt\Delta)/(2a)$ may suffer cancellation. Rationalization gives
\[
 \frac{-b-\sqrt\Delta}{2a}
   =\frac{2c}{-b+\sqrt\Delta},
\]
provided the denominator is nonzero. The analogous rationalization
applies to the other root when $b>0$. Moreover, forming $b^2$ can
overflow even if the roots are finite. For $b\ne0$, the identity
\[
 \sqrt{b^2-4ac}=|b|\sqrt{1-\frac{4ac}{b^2}}
\]
illustrates a rescaling that avoids directly squaring a large $b$;
its implementation must also guard the intermediate quotient.
When $|4ac/b^2|\ll1$, the Taylor approximation
$\sqrt{1-t}=1-\tfrac12t+O(t^2)$ explains the sensitivity of the
small root. In particular, for $b^2\gg|4ac|$,
\[
 x_-\approx\begin{cases}
 -b/a+c/b,&b>0, \\
 -c/b,&b<0.
 \end{cases}
\]
Higher-order terms follow from
$\sqrt{1-t}=1-\tfrac12t-\tfrac18t^2+O(t^3)$.
\end{example}

\begin{example}[Two Ways to Compute a Mean]
For $a, b\in\mathbb F$ compare
\[
 \widehat m=(a\oplus b)\odiv2,
 \qquad \widehat m'=(a\odiv2)\oplus(b\odiv2).
\]
If $|a+b|\le F_{\max}$, the first is correctly rounded:
\[
 \widehat m=\rnd\!\left(\frac{a+b}{2}\right),
 \qquad \min(a, b)\le\widehat m\le\max(a, b).
\]
It may nevertheless overflow when $|a+b|$ exceeds the finite range,
even though the mean is representable. The second avoids such
intermediate overflow; it is correctly rounded when the inputs are
sufficiently large, and otherwise has ULP error at most $1$.
\end{example}

\begin{theorem}[Correctly Rounded Mean]
Let $C\in[2^{e_{\min}+M+3},2^{e_{\max}}]\cap\mathbb F$. For
$a, b\in\mathbb F$, define
\[
 m(a, b)=\begin{cases}
 (a\oplus b)\odiv2,&|a|<C\ \text{and}\ |b|<C, \\
 (a\odiv2)\oplus(b\odiv2),&\text{otherwise}.
 \end{cases}
\]
Then, assuming the stated floating-point operations are defined,
\[
 m(a, b)=\rnd\!\left(\frac{a+b}{2}\right)
 \qquad(a, b\in\mathbb F).
\]
\end{theorem}
